\bookchapter{الفصل 6: احمِ حياتك الرقمية — الأمان والخصوصية}

\chapterintro{أمان الهاتف = أمان الخصوصية + أمان الممتلكات. في هذا العصر الذي أصبحت فيه البيانات أصلاً، حماية هاتفك هي حماية لنفسك.}

\booksection{أولاً: قفل الشاشة: الخط الدفاعي الفيزيائي الأول}

قفل الشاشة هو الخط الدفاعي الأول لهاتفك، وهو الأكثر تجاهلاً.

\booksubsection{ترتيب مستوى الأمان:}

\begin{enumerate}
  \item \textbf{كلمة المرور/رمز PIN}: الأعلى أماناً، وهي الضمان النهائي (عند نسيان البصمة/الوجه تعتمد عليها).
  \item \textbf{التعرف بالبصمة}: التوازن المثالي بين الأمان والراحة.
  \item \textbf{التعرف بالوجه}: الأكثر راحة، لكن الأمان يعتمد على التقنية (التعرف ثلاثي الأبعاد بالهيكل آمن، أما الصور ثنائية الأبعاد فسهلة الخداع).
  \item \textbf{فتح النمط}: مريح لكنه يترك آثاراً يسهل تخمينها.
\end{enumerate}

\booksubsection{ضبط قفل شاشة قوي:}

\begin{itemize}
  \item تأكد من تعيين كلمة مرور معقدة (لا تستخدم تاريخ الميلاد أو 123456).
  \item فعّل التعرف البيومتري.
  \item اضبط وقت قفل الشاشة على"فوري" أو الأقصر.
\end{itemize}

\booksection{ثانياً:“العثور على هاتفي”: الأداة النهائية لمكافحة السرقة}

\textbf{يجب تفعيلها!}يمكنها:

\begin{itemize}
  \item \textbf{التحديد}: معرفة مكان الهاتف على الخريطة.
  \item \textbf{تشغيل الصوت}: جعل الهاتف يرن، مناسب للبحث في المنزل.
  \item \textbf{وضع الفقدان}: قفل الهاتف عن بُعد وعرض رقم اتصال.
  \item \textbf{المسح عن بُعد}: الملاذ الأخير، مسح جميع البيانات (لكن بعد المسح لا يمكن التحديد مرة أخرى).
\end{itemize}

اذهب إلى [الإعدادات]، ابحث عن“Find My iPhone”في Apple، أو “Find My Device” أو خدمة مشابهة من العلامة في Android، وتأكد من أنها مفعّلة باللون الأخضر.

\booksection{ثالثاً: تحديث النظام: ترقيع الثغرات وسد العيوب}

الهدف الأول لتحديث النظام هو إصلاح الثغرات الأمنية، تماماً مثل ترميم ثقوب جدار المنزل. لذلك، عند رؤية تنبيه تحديث النظام، حاول الترقية في أقرب وقت.

بالطبع، قد تحدث مشكلات توافق في التحديثات الكبيرة (مثل من iOS 17 إلى iOS 18)، فيمكنك الانتظار بضعة أيام قبل الترقية.

\booksection{رابعاً: التعرف على الاحتيالعبر الاتصالات ومنعه}

في ظل التطور السريع لتقنيات الذكاء الاصطناعي، فقد نموذج الإدراك التقليدي“رؤية تصدق”فعاليته تماماً. يستغل المحتالون تقنيات الذكاء الاصطناعي مثل التزييف العميق (Deepfake)، لتحقيق ترقية شاملة من الاحتيال“الشبكي الواسع”إلى الجريمة“الاستهداف الدقيق”.

\booksubsection{أكثر أساليب الاحتيال شيوعاً في عصر الذكاء الاصطناعي}

\textbf{(1) احتيال تبديل الوجه وتقليد الصوت بالذكاء الاصطناعي}\par

يجمع المحتالون بطريقة غير مشروعة صوراً عالية الدقة ينشرها المستهدفون على وسائل التواصل الاجتماعي، أو عينات صوتية لا تتجاوز بضع ثوانٍ، ويستخدمون تقنيات التزييف العميق بالذكاء الاصطناعي لاستنساخ“فيديو للقريب” أو “صوت”يكاد يكون مطابقاً. السيناريوهات الشائعة تشمل:\textbf{حوادث مفاجئة، اقتراض طارئ، تعرّض للاختطاف}وغيرها. يستغل هذا الاحتيال الروابط العاطفية كالأخوة والصداقة،\textbf{وضرره الخداعي عالٍ جداً}.

\textbf{(2) انتحال شخصية الجهات الرسمية/خدمة العملاء}\par

تستطيع تقنيات الذكاء الاصطناعي محاكاة أصوات ونبرات موظفي خدمة العملاء الرسميين ومسؤولي الأجهزة الأمنية والقضائية بدقة، بل وتوليف لقطات تحويل مصرفي ووثائق إلكترونية وأوامر توقيف وغيرها من المستندات العالية الواقعية، ومن خلال\textbf{خلق حالة ذعر أو ضغط سلطوي}، يقومون بإغراء الضحايا بالتحويل أو تفعيل مشاركة الشاشة، وبالتالي سرقة معلومات الحساب.

\textbf{(3) احتيال الاستثمار وإدارة الثروات بالذكاء الاصطناعي}\par

عبر خوارزميات الذكاء الاصطناعي تُولَّف لقطات أرباح مزيفة“مكسب مضمون”، أو واجهات برامج اختيار أسهم كمية افتراضية، أو تقارير تحليل استثماري“على مستوى الخبراء”، وبإغراء العائد المرتفع، تُباع خدمات استثمارية باهظة أو يُغرى الضحايا بتحويل أموالهم إلى حسابات تداول يتحكمون بها،\textbf{وفي النهاية يهربون بالأموال}.

\booksubsection{كيفية الكشف — فكرة الحماية بـ“المرآة الكاشفة”}

يمكن للذكاء الاصطناعي استنساخ الصوت والوجه، لكنه ليس منيعاً، وهناك عيوب يمكن استغلالها:

\begin{itemize}
  \item \textbf{طريقة التحقق بالحركة}: عند إجراء مكالمة فيديو، اطلب بشكل استباقي من الطرف الآخر تنفيذ تعليمات حركية متتابعة، مثل"مرّر يدك بسرعة أمام وجهك"، أو "اعرضمحيط الأذن الأذن جانبياً"، أو "أغمض عينيك بسرعة 5 مرات". تقنيات تبديل الوجه الحالية بالذكاء الاصطناعي عند معالجة الحجب الديناميكي المعقد والمشاهد سريعة الحركة وتغيرات الضوء والظل،\textbf{تظهر بسهولة عيوب مثل تجميد الصورة وتشوه الوجه أو اختراق الحواف}.
  \item \textbf{طريقة التحقق بالشفرات}: أنشئ واستخدم أسئلة تحقق خاصة مسبقة الإعداد، مثل"ما الطبق الذي طلبته في آخر مرة تناولنا العشاء معاً"،"ما أول هدية أهديتك إياها"وغيرها من التفاصيل التي لا يعرفها إلا الطرفان. إذا تردد الطرف الآخر بشكل ملحوظ أو أجاب خارج الموضوع، يجب\textbf{إنهاء المكالمة فوراً وتفعيل خطة الطوارئ الأمنية}.
  \item \textbf{طريقة التحقق الثانوي}: التزم بشكل صارم بمبدأ"\textbf{أربعة لا وتحقق واحد}"الذي أصدره مكتب الأمن السيبراني بوزارة الأمن العام — لا تصدق المعلومات الغريبة، ولا تحوّل الأموال بشكل عشوائي، لا تكشف المعلومات الشخصية الحساسة، لا تنشر المحتوى المشبوه. عند أي تعامل مالي، تأكد من خلال\textbf{إعادة الاتصال بالرقم الأصلي المعروف أو التأكيد وجهاً لوجه}. عند إجراء التحويل، يُستحسن اختيار وضع"الوصول خلال 24 ساعة"، لترك نافذة زمنية كافية لتحديد الاحتيال المحتمل واسترداد الأموال.
\end{itemize}

\booksubsection{الأساليب التقليدية الشائعة (لا يزال ينبغي الحذر منها):}

\begin{itemize}
  \item \textbf{الفوز بالجائزة/استرداد العمولة}: إغراء بالمنفعة، يمنحك شيئاً صغيراً أولاً ثم يخدعك بمبلغ كبير.
  \item \textbf{احتيال العلاقات عبر الإنترنت}: تسمى"حظيرة الخنازير"، تنمية المشاعر ثم الاحتيال المالي.
  \item \textbf{انتحال شخصية المعارف}: استعارة الأموال بعد سرقة الحساب (أصعب تمييزاً في عصر الذكاء الاصطناعي).
\end{itemize}

\booksubsection{الحماية المدمجة في الهاتف:}

\begin{itemize}
  \item حظر المكالمات المزعجة.
  \item منع سرقة رمز التحقق.
  \item الفحص الأمني لروابط الرسائل النصية.
\end{itemize}

\booksubsection{نشيد مكافحة الخداع:}

\textbf{لا تصدق، لا تكشف، لا تنقر، لا تحوّل، وعند الشك تحقّق أولاً!}عند مواجهة“أمر طارئ”ابقَ هادئاً 5 دقائق، فكل عمليات الاحتيال تقريباً تنكشف في هذه اللحظة.

\booksection{خامساً: مقاومة“الفتح”وهجمات الهندسة الاجتماعية: قطع سلسلة تسرب الخصوصية}

المبدأ الأساسي لـ“الفتح”(هجمات الهندسة الاجتماعية) هو\textbf{الجمع والتحليل المنهجي للشظايا الرقمية الشخصية وبناء صورة شخصية كاملة}. مفتاح الوقاية هو رسم خطوط حمراء واضحة للخصوصية الشخصية، وبناء نظام حماية متعدد الطبقات.

\booksubsection{معالجة إخفاء بيانات الصور}

طرق التنظيف الأساسية لـ EXIF موجودة في الفصل 2 “إدارة الصور ونظام الملفات”. قبل مشاركة الصور على وسائل التواصل الاجتماعي، يجب اتخاذ خطوات إضافية\textbf{لحذف بيانات الصور الوصفية (معلومات EXIF) بالكامل}، بما في ذلك وقت التصوير وموديل الجهاز وإحداثيات موقع GPS وغيرها من البيانات الحساسة. يجب أيضاً تشفير أو طمس المعلومات الرئيسية في الصور مثل\textbf{أرقام المنازل، ومعلومات لوحات السيارات، والرموز الشريطية للتذاكر، وأرقام بطاقات الهوية}، لتجنب تسريب مسار تحركاتك الشخصية دون قصد.

\begin{itemize}
  \item iOS: يمكن إغلاق"الموقع الدقيقة"في [الإعدادات] > [الخصوصية والأمان] > [الكاميرا]؛ وعند المشاركة، اختر"الصورة الحالية فقط"لإزالة البيانات الوصفية.
  \item Android: تطبيقات معرض الصور (مثل Google Photos) تسأل عند المشاركة ما إذا كنت تريد الاحتفاظ بمعلومات الموقع.
  \item أدوات عامة: ExifTool وMetapho وغيرها لعرض/تحرير EXIF.
\end{itemize}

\booksubsection{تقليص أذونات التطبيقات}

طرق إدارة الأذونات الأساسية موضحة بالتفصيل في الفصل 4 “إدارة الأذونات”. في سيناريو “منع الكشف عن الهوية”، يجب مراجعة وإغلاق أذونات تطبيقات التجارة الإلكترونية والأدوات بشكل دوري\textbf{الميكروفون وجهات الاتصال ومعلومات الموقع}، لمنع التطبيقات من“التنصت”على المحادثات اليومية في الخلفية أو جمع شبكة العلاقات الشخصية بشكل غير قانوني؛ يجب إلغاء تثبيت التطبيقات المهجورة التي لا تزال تشغل أذونات عالية، لتقليل نقاط تسرب الخصوصية من المصدر.

\begin{itemize}
  \item عند رؤية طلب"إذن جهات الاتصال"، اسأل أولاً: هل يحتاج هذا التطبيق فعلاً إلى معرفة جميع جهات اتصالي؟
  \item التطبيقات التي لم تستخدمها منذ ثلاثة أشهر، ألغِ تثبيتها بكل حزم. الأذونات تُشدَّد مع الاستخدام، لا تُرخى.
\end{itemize}

\booksubsection{تأمين الأجهزة الذكية}

عند استخدام الكاميرات الذكية وأقفال الأبواب الذكية ومكبرات الصوت الذكية وغيرها من أجهزة IoT لأول مرة، يجب\textbf{تغيير كلمة مرور الدخول الافتراضية فوراً}، وتعيين كلمة مرور قوية مستقلة تحتوي على أحرف كبيرة وصغيرة وأرقام ورموز خاصة؛ يجب إبقاء وظائف الوصول عن بُعد غير المستخدمة بشكل دائم مغلقة، وفحص تحديثات البرامج الثابتة للأجهزة بانتظام وترقيع الثغرات الأمنية في الوقت المناسب،\textbf{لمنع اختراق القراصنة والتجسس على مساحة الخصوصية المنزلية}.

\begin{tcbwarn}{تحذير من حالة واقعية}
في 2022، أبلغت الشرطة في منطقة ما أن القراصنة من خلال كسر كلمات المرور الافتراضية لكاميرات منزلية من علامة معينة (لم يغيّرها كثير من المستخدمين أبداً)، سرقوا لقطات غرف نوم آلاف الأسر وعرضوها للبيع عبر الإنترنت.\textbf{عدم تغيير كلمة المرور الافتراضية = ترك باب منزلك مفتوحاً}.

\end{tcbwarn}

\booksection{سادساً: حماية الخصوصية: بياناتك هي ممتلكاتك}

الصور وسجلات الدردشة ومسارات الموقع في هاتفك هي في جوهرها“أصولك الرقمية”. بمجرد تسريبها، قد يؤدي ذلك إلى إزعاج إعلاني في أحسن الأحوال، أو إلى التشهير والاحتيال وحتى الابتزاز في أسوأها.

\booksubsection{مبدأ“الحد الأدنى”لأذونات التطبيقات:}

\begin{itemize}
  \item \textbf{الموقع}: تطبيقات استدعاء السيارة/الخرائط تستخدم"السماح دائماً"؛ التطبيقات الأخرى تستخدم"السماح أثناء الاستخدام فقط".
  \item \textbf{جهات الاتصال}: لا تمنح إلا عند الضرورة (مثل إضافة الأصدقاء في WeChat)، وإلا فامنعها دائماً.
  \item \textbf{معرض الصور/جميع الصور}: من iOS 14+ يمكن منح"صور مختارة"، ويدعم Android 13+ أيضاً التفويض الدقيق.
  \item \textbf{الميكروفون/الكاميرا}: أغلقها عادة، وافتحها عند الحاجة. عند الاشتباه في التسجيل الخفي، اسحب مركز التحكم لرؤية ما إذا كانت هناك"نقطة خضراء صغيرة/نقطة برتقالية صغيرة".
  \item \textbf{التتبع في الخلفية}: من iOS 14.5+ يُلزم التطبيقات بالسؤال قبل استخدام IDFA؛ يمكن إغلاق"معرّف الإعلان"في Android.
\end{itemize}

\booksubsection{تطبيق عملي لإدارة الخصوصية:}

\begin{itemize}
  \item \textbf{تقرير الخصوصية (iOS)}: [الإعدادات] > [الخصوصية والأمان] > [تقرير خصوصية التطبيقات]، لعرض التطبيقات التي تصل خلسة إلى موقعك وكاميرتك وميكروفونك.
  \item \textbf{سجل استخدام الأذونات (Android)}: [الإعدادات] > [الخصوصية] > [سجل استخدام الأذونات]، لمعرفة من استدعى الأذونات الحساسة ومتى.
  \item \textbf{تقييد التتبع الإعلاني}: أغلق"معرّف الإعلان" أو "تقييد التتبع الإعلاني"في إعدادات النظام، لتقليل"الإزعاج بالإعلانات المستهدفة".
\end{itemize}

\booksection{سابعاً: قائمة“الإسعافات الأولية”بعد الفقد/السرقة}

لا تذعر عند فقدان الهاتف، تعامل بالترتيب التالي:

\begin{enumerate}
  \item \textbf{تجميد بطاقة SIM فوراً}: اتصل بالمشغل للإبلاغ عن الفقدان، لمنع استخدامها في استقبال رموز التحقق.
  \item \textbf{القفل عن بُعد}: فعّل"وضع الفقدان"عبر الخدمات السحابية (iCloud/سحابة العلامة).
  \item \textbf{تغيير كلمات المرور المهمة}: غيّر فوراً كلمات المرور لجميع الحسابات المهمة التي سجّلت الدخول إليها من الهاتف (الدفع، البنوك، البريد الإلكتروني).
  \item \textbf{الإبلاغ}: أبلغ الشرطة، وسجّل رقم IMEI. يسهّل ذلك إثبات الهوية عند استرجاع الهاتف.
  \item \textbf{إخطار الأهل والأصدقاء}: لمنع احتيال المحتالين باسمك لاقتراض الأموال.
  \item \textbf{(أخيراً) المسح عن بُعد}: نفّذ هذه العملية فقط عند التأكد من استحالة الاسترجاع، حيث تفقد وظيفة التحديد بعد المسح.
\end{enumerate}

\booksection{ثامناً: أمان بطاقة SIM وeSIM}

\begin{itemize}
  \item \textbf{تعيين PIN لبطاقة SIM}: بعد التفعيل، يتطلب كل إعادة تشغيل للهاتف أو تغيير البطاقة إدخال رمز PIN، لمنع تجاوز قفل الشاشة عبر"سحب البطاقة وتبديل الهاتف".
  \item \textbf{مزايا eSIM}: لا توجد بطاقة مادية يمكن سحبها، أمان أعلى؛ كما تدعم التبديل بين أرقام متعددة، مريحة جداً أثناء السفر.
  \item \textbf{مخاطر تبديل البطاقة الثانوي}: احذر من احتيال"تبديل البطاقة الثانوي للمشغل"— يستخدم المحتالون هويتك لإصدار بطاقة SIM بديلة، وبالتالي السيطرة على رقم هاتفك. تعيين"إصدار بطاقة بديلة يتطلب حضور صاحب الخط بشهادة الهوية إلى فرع المشغل"هو الخط الدفاعي الأكثر فعالية.
\end{itemize}

\begin{tcbwarn}{أشياء يجب ألا تفعلها أبداً}
\begin{itemize}
  \item استخدام كلمات مرور ضعيفة مثل تاريخ الميلاد/123456/qwerty؛
  \item إرسال رمز التحقق لأي شخص (حتى لو كان"خدمة العملاء")؛
  \item النقر على روابط غير معروفة في الرسائل النصية، خاصة من الأرقام الأجنبية التي تبدأ بـ 00؛
  \item تسجيل الدخول إلى البنوك الإلكترونية أو الدفع تحت Wi-Fi عام؛
  \item تفعيل"مشاركة الشاشة" أو "التحكم عن بُعد"للغرباء.
\end{itemize}

\end{tcbwarn}

\begin{tcbpractice}{جرّب تفعيل ومحاكاة"القفل عن بُعد"}
\begin{enumerate}
  \item تأكد من تفعيل وظيفة"العثور على هاتفي".
  \item على متصفح كمبيوتر أو هاتف آخر، سجّل الدخول إلى حساب الخدمة السحابية المقابل (iCloud.com أو موقع الخدمة السحابية للعلامة).
  \item ابحث عن وظيفة"العثور على الجهاز"، وحاول إطلاق أمر"تشغيل الصوت" أو محاكاة أمر"القفل"على هاتفك الخاص، للتعرف على العملية بالكامل.
\end{enumerate}

\end{tcbpractice}
